\section{Discussion}\label{sec:discussion}

\subsection{Limitations}\label{sec:limitations}

\paragraph{Activity Realism.}
The mean JS divergence of 0.218 indicates moderate distributional
similarity between \sys{}-generated and real-world activity
patterns.  Divergence is strongly dataset-dependent: ARAS datasets
(ground-truth labels, 27~activity types) show low JS (0.101--0.113),
while CASAS datasets (room-proxy labels, 3--4~categories) show
higher JS (0.278--0.301) due to category-coverage mismatch.
Notably, \sys{} achieves strong temporal correlation on CASAS
(0.80--0.87), confirming realistic circadian patterns.
A leave-one-out Markov-chain baseline achieves lower JS on CASAS
(which favors data-trained models on homogeneous datasets) but
\emph{higher} JS on ARAS (0.202--0.297 vs.\ 0.101--0.113),
demonstrating \sys{}'s advantage on diverse activity profiles
without requiring training data.
The CASAS room-proxy labeling remains a limitation:
since the publicly available CASAS sensor logs
lack explicit activity annotations, we infer categories from sensor
locations (e.g., \emph{Kitchen}$\to$\emph{Meal\_Preparation}).
This proxy is deterministic and consistent but conflates
multi-purpose rooms, inflating divergence metrics.

\paragraph{Sim2Real Transfer.}
With enriched features (transition matrices, n-grams, entropy),
transfer accuracy improves substantially over our initial coarse
features: ARAS-House~A achieves $0.813$~RF accuracy
($F_1 = 0.783$), up from $0.075$ with hourly distributions alone.
However, domain discrimination remains near~$1.0$ on CASAS datasets,
indicating that \sys{}'s 9-category schedules are trivially separable
from CASAS room-proxy data.  Practical Sim2Real transfer benefits
from datasets with rich, diverse activity annotations;
researchers targeting CASAS-style deployments would benefit from
domain adaptation with a small number of real labeled samples.

\paragraph{Navigation Policy.}
The SPL of~1.0 follows from our use of a shortest-path follower
on the navmesh.  This validates the goal reachability pipeline
but does not evaluate learned navigation policies, which is
orthogonal to our platform contribution.

\subsection{Threats to Validity}

\emph{Internal.}
LLM-generated schedules are non-deterministic even with fixed seeds
due to floating-point non-determinism in quantized inference.  We mitigate
this by generating multiple schedules per configuration and reporting
confidence intervals.  In the autonomous evaluation, GPT-OSS~20B
produced 4{,}307~scheduled tasks across 28~scenes, with per-scene
task counts ranging from 3 to 190.

\emph{External.}
Our reference datasets (CASAS, ARAS) represent Western households
with sensor-based annotations; cultural and geographic diversity is
limited.  The persona library partially addresses this but cannot
fully capture global variation.  The 28~HSSD scenes, while diverse
(6--35~rooms), represent a subset of possible residential layouts.

\emph{Construct.}
We use KL/JS divergence and temporal correlation as proxies for
``realism.''  Human evaluation studies would complement these
automated metrics but are beyond the scope of this work.

\paragraph{Firmware Fidelity.}
QEMU's ARM emulation faithfully executes instruction-level code but
does not model analog peripherals (e.g., precise ADC noise profiles),
RF propagation, or cycle-accurate timing.  For protocols above the
physical layer (UART, MQTT, HTTP), the emulation is functionally
exact.  Our autonomous evaluation
processed 47{,}207~firmware state changes through the full pipeline
(GPIO interrupt $\to$ ISR $\to$ state update $\to$ SmartThings cloud)
with zero data loss, confirming protocol-level fidelity at scale.

\paragraph{LLM Limitations.}
All schedule generation uses GPT-OSS~20B, a 20-billion-parameter model
served locally via LMStudio with 4-bit quantization.
Our validation layer catches JSON parse errors and retries with explicit
error feedback, bringing the effective failure rate below~1\%.
The retry mechanism with increased temperature (0.9) proved effective:
across the 28-scene autonomous evaluation, the model generated
4{,}307~tasks with only 2~scenes failing to navigate (due to navmesh
issues, not LLM errors).
Future work could evaluate additional LLM models to assess the
impact of model choice on schedule diversity and correctness.

\paragraph{Scalability Ceiling.}
While \sys{} handles 200~event-bus devices efficiently (10{,}901~events/s,
35\% CPU), Docker QEMU containers introduce per-container overhead.
The autonomous evaluation ran 6~containers per scene across
28~scenes (88.0~hours total) without resource exhaustion; memory
grew by at most 24.7\,MB over 168~simulated hours.  Lightweight
container runtimes (e.g., gVisor) or multi-machine orchestration
(Kubernetes) could raise this ceiling further.

\paragraph{Navigation Reliability.}
The 94.9\% navigation success rate (3{,}580~trials across 26~navigable
scenes) demonstrates that the combination of navmesh snapping and
reachability filtering produces reliable goal positions.  The 5.1\%
failure rate consists of edge cases where humanoid--furniture
collisions or multi-floor navmesh discontinuities caused
goal-reaching timeouts.  Two additional scenes (102816009 and
106366371\_174226743) failed to produce any navigation trials due
to navmesh generation failures, reducing the navigable scene count
from 28 to 26.  Room coverage averages 62.3\% (95\% CI: [55.3\%,
69.3\%]), with upper floors and isolated outdoor areas accounting
for the majority of unreachable rooms.

\paragraph{Privacy and Ethics.}
\sys{} generates synthetic behavioral data and does not collect or
process real occupant data.  Personas are fictional.  The SmartThings
integration uses the researcher's own account; no third-party
accounts are accessed.

\paragraph{Security Research Platform.}
The successful reproduction of the phantom-delay
attack~\cite{phantom-delay} within \sys{}
(\Cref{sec:eval-phantom-delay}), combined with the addition of
two new attack suites---Malicious SmartApp (Suite~4, CVSS~8.8)
targeting cloud API authorization flaws, and ESP32 Buffer Overflow
(Suite~5, CVSS~9.8) targeting firmware memory safety---demonstrates
a broader capability: \sys{} can serve as a \emph{reproduction and
discovery platform} for IoT security research.
The SmartApp attack exposed that all 5 registered devices were
enumerable without valid OAuth credentials, and a forged DISARM
command was accepted in 5\,ms.  The ESP32 overflow caused device
crash/DoS with a 137-byte payload, with recovery after $\sim$12\,s.
Crucially, all attack traffic is captured as genuine
\texttt{.pcap} files via \texttt{tshark}
(\Cref{sec:eval-pcap}): the 154{,}151-segment global capture
and 736~per-attack pcaps provide independently verifiable
evidence that attacks traverse the real Docker bridge
network---not an in-process simulation.
Researchers can implement and evaluate protocol-level, cloud API,
and memory-safety exploits within the 3D simulation without
requiring physical hardware, reducing the cost and complexity of
IoT security research.

\paragraph{Broader Impact.}
By providing a high-fidelity simulation environment, \sys{} reduces
the need for real-world deployments during early-stage IoT research,
lowering costs and privacy risks.  However, realistic synthetic
behavioral data could be misused to train surveillance systems;
we encourage responsible use aligned with ethical guidelines.
